\subsection{Checkpoint 调试清单}

Checkpoint 往往不是“能不能存”这么简单，而是“恢复后是不是同一个实验”。如果恢复后的 loss 曲线明显不一致，通常优先查下面几项：

\begin{table}[H]
\centering
\begin{tabular}{p{0.22\textwidth}p{0.28\textwidth}p{0.26\textwidth}}
\toprule
\textbf{现象} & \textbf{优先检查} & \textbf{常见原因} \\
\midrule
恢复后 loss 跳变 & optimizer state、学习率调度 & 只存了 model 没存 optimizer \\
训练变慢 & device / dtype / data loader & batch 没放到 GPU 或 pinned memory \\
梯度不更新 & requires\_grad / zero\_grad & 参数被冻住或梯度没清 \\
结果不可复现 & RNG、数据顺序、dropout & 随机种子没统一 \\
\bottomrule
\end{tabular}
\caption{checkpoint 的问题通常要从训练状态而不是单纯参数开始排查}
\end{table}

